\chapter{Conclusiones}\label{cap:conclusiones}

\section{Síntesis del trabajo realizado}
Este proyecto ha desarrollado una solución analítica completa para estudiar movilidad
urbana con datos masivos de taxis de Nueva York, combinando Spark, Zeppelin y
Python/PySpark en un entorno local reproducible. El resultado no se limita a un
notebook aislado: se ha construido un flujo técnico coherente que integra ingesta
de Parquet, calidad del dato, análisis tabular y gráfico, exportación de indicadores
por zona y visualización geográfica con mapas interactivos.

\section{Cumplimiento de objetivos}
Los objetivos definidos en la introducción se consideran alcanzados:
\begin{itemize}
	\item se cargó y validó el dataset de enero 2025 en formato Parquet;
	\item se implementó limpieza trazable con criterios explícitos y defendibles;
	\item se generaron variables derivadas para análisis temporal y económico;
	\item se ejecutaron consultas analíticas en Spark SQL y DataFrames;
	\item se enriquecieron los datos geográficamente con GeoJSON de zonas TLC;
	\item se exportaron 14 CSVs de indicadores por zona de origen y destino;
	\item se construyeron dos notebooks de Zeppelin: uno con análisis completo
	y otro con mapas coropléticos interactivos con Folium.
\end{itemize}

\section{Valor técnico y académico}
El valor principal del trabajo reside en su equilibrio entre ingeniería y
comunicación de resultados:
\begin{itemize}
	\item desde el punto de vista técnico, el flujo es reproducible con Docker y
	fácilmente extensible a más meses o datasets;
	\item desde el punto de vista académico, los notebooks permiten defender
	decisiones de diseño con evidencia ejecutable;
	\item desde el punto de vista metodológico, se mantuvo trazabilidad entre
	requisitos, implementación y pruebas;
	\item la separación en dos notebooks (análisis + mapas) facilita la
	reutilización y la presentación modular del trabajo.
\end{itemize}

\section{Limitaciones}
Se reconocen limitaciones inherentes al alcance elegido:
\begin{itemize}
	\item infraestructura local con recursos acotados;
	\item ausencia de despliegue productivo con alta disponibilidad;
	\item análisis limitado a un único mes (enero 2025), que puede no
	reflejar patrones estacionales completos.
\end{itemize}

Estas limitaciones son compatibles con el objetivo del proyecto y no comprometen la
validez de las conclusiones dentro del contexto definido.

\section{Líneas de trabajo futuro}
Como evolución natural del proyecto se proponen:
\begin{itemize}
	\item ampliar el periodo temporal analizado con más meses y temporadas completas;
	\item desplegar el flujo en un clúster real (Standalone, YARN o Kubernetes);
	\item incorporar orquestación de pipelines y pruebas automáticas de regresión;
	\item añadir capas de análisis predictivo (por ejemplo, demanda por franja horaria);
	\item construir un panel interactivo operativo para consumo no técnico;
	\item incluir una comparativa de rendimiento formal entre Parquet y CSV con
	infraestructura controlada.
\end{itemize}

\section{Cierre}
La solución final demuestra que es posible construir, con herramientas abiertas y
en un entorno docente, un sistema de analítica Big Data robusto y defendible.
Los notebooks de Zeppelin actúan como artefactos de síntesis del proyecto y dejan
una base sólida para continuar hacia escenarios de mayor escala y complejidad.

